
\begin{resumen}
\small

La epilepsia es un trastorno neurológico que requiere una
evaluación clínica especializada. La electroencefalografía
(EEG) desempeña un papel importante en la evaluación de
alteraciones de la actividad eléctrica cerebral. Sin
embargo, la interpretación de los registros
electroencefalográficos puede demandar un tiempo
considerable, especialmente cuando se analizan estudios
de larga duración y grandes volúmenes de datos.

Este trabajo presenta el diseño y la implementación de
EpilexIA, una plataforma web desarrollada para apoyar
a los profesionales de la salud en la gestión,
visualización y análisis de estudios
electroencefalográficos. La plataforma incorpora un
componente de inteligencia artificial diseñado para
apoyar la identificación y revisión de posibles
anomalías en señales EEG.

El sistema se desarrolló utilizando React para la
interfaz web, Java y Spring Boot para los servicios
del \textit{backend} y PostgreSQL para la persistencia
de información. El subsistema de inteligencia
artificial utiliza componentes basados en Python,
FastAPI para la integración de servicios y bibliotecas
científicas para el procesamiento de señales EEG.
La arquitectura separa la gestión de información
clínica de los componentes responsables del análisis
automatizado de señales.

EpilexIA proporciona funcionalidades para la gestión
de pacientes y estudios EEG, visualización de señales
en formatos EDF y CSV, aplicación de filtros, ajuste
de sensibilidad, selección de montajes
electroencefalográficos, anotación de anomalías y
elaboración de informes médicos. La plataforma también
incorpora mecanismos de control de acceso y funciones
para compartir estudios entre profesionales de la
salud autorizados.

La evaluación comprendió pruebas funcionales y de
integración de los principales componentes de gestión
clínica y visualización de señales EEG, así como una
evaluación contextual realizada por una médica
especialista en el Hospital Central del Instituto
de Previsión Social (IPS). Los resultados
proporcionaron evidencia del correcto funcionamiento
de las principales funcionalidades evaluadas.
No obstante, la ejecución completa y el desempeño
clínico del subsistema de inteligencia artificial
no pudieron verificarse íntegramente dentro del
alcance de las pruebas realizadas.

En conclusión, EpilexIA constituye una plataforma
tecnológica de apoyo para la organización,
visualización y análisis de estudios EEG, y establece
una base para futuras actividades de validación y
mejora de la detección de anomalías asistida por
inteligencia artificial.

\end{resumen}